\section{Introduction} \label{sec:introduction}
\textit{Property-based testing}
is an approach to testing programs.
In this approach,
a programmer identifies
\textit{properties}, or specifications,
of functions or groups of functions,
and those properties are then tested on
a large number of randomly selected values.
Property-based testing was first popularized in
the functional programming ecosystem,
with the Haskell library QuickCheck~\cite{claessen2000quickcheck}
being a particular influential work.
However, there are now property-based testing libraries
for many other languages, both imperative and functional,
including libraries for Python~\cite{maciver2019hypothesis},
Rust~\cite{derbenwick2025property},
\bill{C++? Java? 2 or 3 more here maybe}

A traditional style of
property based test
ensures that two expressions return
the same result when applied to equal inputs.
For example
\bill{Show something in quickcheck}
\bill{While undoubtedly useful,
... why might we prefer \ourDSL?
(1) \ourDSL works as a specification
(2) \ourDSL is independent of existing code-- I'm thinking one good example
here could be that zip/zipWith have similar behavior- but also these might have very similar implementations}
In this work,
we propose a
new embedded Domain Specific Language (eDSL),
\ourDSL,
for the purpose of writing property-based tests.